% Options for packages loaded elsewhere
\PassOptionsToPackage{unicode}{hyperref}
\PassOptionsToPackage{hyphens}{url}
\PassOptionsToPackage{table}{xcolor}
%
\documentclass[
  11pt,
]{report}
\usepackage{amsmath,amssymb}
\usepackage{iftex}
\ifPDFTeX
  \usepackage[T1]{fontenc}
  \usepackage[utf8]{inputenc}
  \usepackage{textcomp} % provide euro and other symbols
\else % if luatex or xetex
  % pdfLaTeX-only repository policy: the original unicode-math/fontspec
  % branch is unreachable under pdfLaTeX and is replaced with a guard.
  \PackageError{cpython-suite}{Compile with pdfLaTeX only}{}
\fi
\usepackage{lmodern}
\ifPDFTeX\else
  % xetex/luatex font selection
\fi
% Use upquote if available, for straight quotes in verbatim environments
\IfFileExists{upquote.sty}{\usepackage{upquote}}{}
\IfFileExists{microtype.sty}{% use microtype if available
  \usepackage[]{microtype}
  \UseMicrotypeSet[protrusion]{basicmath} % disable protrusion for tt fonts
}{}
\makeatletter
\@ifundefined{KOMAClassName}{% if non-KOMA class
  \IfFileExists{parskip.sty}{%
    \usepackage{parskip}
  }{% else
    \setlength{\parindent}{0pt}
    \setlength{\parskip}{6pt plus 2pt minus 1pt}}
}{% if KOMA class
  \KOMAoptions{parskip=half}}
\makeatother
\usepackage{xcolor}
\usepackage{listings}
\newcommand{\passthrough}[1]{#1}
\lstset{defaultdialect=[5.3]Lua}
\lstset{defaultdialect=[x86masm]Assembler}
\usepackage{longtable,booktabs,array}
\usepackage{calc} % for calculating minipage widths
% Correct order of tables after \paragraph or \subparagraph
\usepackage{etoolbox}
\makeatletter
\patchcmd\longtable{\par}{\if@noskipsec\mbox{}\fi\par}{}{}
\makeatother
% Allow footnotes in longtable head/foot
\IfFileExists{footnotehyper.sty}{\usepackage{footnotehyper}}{\usepackage{footnote}}
\makesavenoteenv{longtable}
\setlength{\emergencystretch}{3em} % prevent overfull lines
\providecommand{\tightlist}{%
  \setlength{\itemsep}{0pt}\setlength{\parskip}{0pt}}
\setcounter{secnumdepth}{5}
% CPython Engineering Suite shared presentation layer -- pdfLaTeX compatible
\usepackage[T1]{fontenc}
\usepackage{lmodern}
\usepackage{microtype}
\usepackage{geometry}
\geometry{letterpaper,margin=0.91in,headheight=16pt,headsep=19pt}
\usepackage{xcolor} % [table] passed at top of file (option-clash fix)
\definecolor{CPNavy}{HTML}{13283B}
\definecolor{CPBlue}{HTML}{176BA0}
\definecolor{CPGold}{HTML}{D8A73C}
\definecolor{CPMuted}{HTML}{566979}
\definecolor{CPBg}{HTML}{EFF5F8}
\usepackage{titlesec}
\titleformat{\chapter}[display]{\normalfont\sffamily\bfseries\color{CPNavy}}{\large\color{CPBlue}CHAPTER \thechapter}{9pt}{\Huge}[\vspace{6pt}\color{CPGold}\titlerule]
\titleformat{\section}{\Large\bfseries\sffamily\color{CPBlue}}{\thesection}{0.7em}{}
\titleformat{\subsection}{\large\bfseries\sffamily\color{CPNavy}}{\thesubsection}{0.7em}{}
\usepackage{fancyhdr}
\pagestyle{fancy}
\fancyhf{}
\fancyhead[L]{\footnotesize\color{CPMuted}\sffamily CPYTHON ENGINEERING SUITE}
\fancyhead[R]{\footnotesize\color{CPMuted}\sffamily 3.16.0a1+dev SNAPSHOT}
\fancyfoot[C]{\footnotesize\color{CPMuted}\thepage}
\renewcommand{\headrulewidth}{0.3pt}
\AtBeginDocument{\urlstyle{tt}}
\lstset{basicstyle=\ttfamily\small,columns=fullflexible,keepspaces=true,breaklines=true,breakatwhitespace=false,showstringspaces=false,frame=single,framesep=5pt,rulecolor=\color{CPBg},backgroundcolor=\color{CPBg},xleftmargin=3pt,xrightmargin=3pt}
\usepackage{tcolorbox}
\tcbuselibrary{breakable,skins}
\newtcolorbox{sourcebox}[1][]{enhanced,breakable,sharp corners,boxrule=0.4pt,colback=CPBg,colframe=CPBlue,borderline west={3pt}{0pt}{CPBlue},left=10pt,right=8pt,top=6pt,bottom=6pt,#1}
\newtcolorbox{cautionbox}[1][]{enhanced,breakable,sharp corners,boxrule=0.4pt,colback=CPBg,colframe=CPGold,borderline west={3pt}{0pt}{CPGold},left=10pt,right=8pt,top=6pt,bottom=6pt,#1}
\usepackage{tikz}
\usetikzlibrary{arrows.meta,positioning,shapes.geometric,calc}
\usepackage{enumitem}
\setlist{itemsep=2pt,topsep=4pt}
\usepackage{booktabs}
\usepackage{longtable}
\AtBeginEnvironment{longtable}{\small}
\usepackage{array}
\setlength{\parskip}{5pt}
\setlength{\parindent}{0pt}
\setcounter{tocdepth}{2}
\usepackage{hyperref}
\hypersetup{colorlinks=true,linkcolor=CPBlue,urlcolor=CPBlue,citecolor=CPBlue,pdfauthor={CPython Engineering Documentation Suite}}
\makeatletter
\renewcommand{\maketitle}{%
  \begin{titlepage}\pagecolor{CPNavy}\color{white}\sffamily
  \vspace*{0.8in}
  {\color{CPGold}\rule{2.1in}{4pt}}\\[16pt]
  {\Large\bfseries CPYTHON ENGINEERING DOCUMENTATION SUITE}\\[24pt]
  {\Huge\bfseries \@title\par}\vspace{19pt}
  \vfill
  {\normalsize Source baseline: CPython 3.16.0 alpha 1+dev repository snapshot}\par
  {\normalsize Prepared 16 October 2026}\par
  {\small\color{white!70} Source-traced engineering commentary; not an official Python specification.}\vspace{0.65in}
  \end{titlepage}\nopagecolor\color{black}
  }
\makeatother
\ifLuaTeX
  \usepackage{selnolig}  % disable illegal ligatures
\fi
\usepackage{bookmark}
\IfFileExists{xurl.sty}{\usepackage{xurl}}{} % add URL line breaks if available
\urlstyle{same}
\hypersetup{
  pdftitle={Python Testing, Security and Performance Handbook},
  hidelinks,
  pdfcreator={LaTeX via pandoc}}

\title{Python Testing, Security and Performance Handbook}
\usepackage{etoolbox}
\makeatletter
\providecommand{\subtitle}[1]{% add subtitle to \maketitle
  \apptocmd{\@title}{\par {\large #1 \par}}{}{}
}
\makeatother
\subtitle{Volume 07 \textbar{} Verification strategy, AppSec controls,
diagnostics and CI}
\author{}
\date{}

\begin{document}
\maketitle

{
\setcounter{tocdepth}{1}
\tableofcontents
}
\chapter{Engineering assurance model}\label{engineering-assurance-model}

\begin{sourcebox}
\textbf{Evidence classification.} \textbf{SOURCE} = statements supported by material present in the uploaded snapshot; \textbf{IMPLEMENTATION} = CPython-specific internals, not guaranteed Python semantics; \textbf{ENGINEERING} = original recommendations, examples, checks, and exercises. The source baseline is \textbf{3.16.0 alpha 1+dev}; missing repository file bodies are noted as evidence limitations.
\end{sourcebox}

This handbook combines CPython's source, docs and testing/tooling
inventory with original engineering controls. The repository shows
\passthrough{\lstinline!Lib/test/!},
\passthrough{\lstinline!Doc/library/!} and
\passthrough{\lstinline!Tools/!}; its internal architecture
documentation points contributors to specialized compiler, runtime and
extension behavior. \textbf{Security guidance and CI examples here are
authored recommendations, not normative CPython policy.} The uploaded
snapshot represents 3.16.0 alpha 1+dev; every production control must be
validated against its actual interpreter version and dependencies.

\section{Assurance dimensions}\label{assurance-dimensions}

\begin{longtable}[]{@{}
  >{\raggedright\arraybackslash}p{(\columnwidth - 4\tabcolsep) * \real{0.3333}}
  >{\raggedright\arraybackslash}p{(\columnwidth - 4\tabcolsep) * \real{0.3333}}
  >{\raggedright\arraybackslash}p{(\columnwidth - 4\tabcolsep) * \real{0.3333}}@{}}
\toprule\noalign{}
\begin{minipage}[b]{\linewidth}\raggedright
Dimension
\end{minipage} & \begin{minipage}[b]{\linewidth}\raggedright
Goal
\end{minipage} & \begin{minipage}[b]{\linewidth}\raggedright
Evidence artifact
\end{minipage} \\
\midrule\noalign{}
\endhead
\bottomrule\noalign{}
\endlastfoot
Correctness & Specified inputs produce expected outputs/errors & Unit,
integration and property tests \\
Robustness & Failures are explicit and contained & Fault injection and
cleanup tests \\
Security & Trust boundaries preserve confidentiality and integrity &
Abuse cases and security tests \\
Reliability & System recovers from crashes/timeouts/cancellation &
End-to-end restart tests \\
Performance & Capacity targets met without invalid assumptions &
Repeatable benchmark and profile \\
Maintainability & Behavior is observable and change-controlled & Source
map, CI results and review decisions \\
\end{longtable}

\chapter{Testing structure}\label{testing-structure}

\section{Unit, integration and system
tests}\label{unit-integration-and-system-tests}

\passthrough{\lstinline!unittest!},
\path{unittest.mock},
\passthrough{\lstinline!doctest!}, and the CPython test suite illustrate
layered verification. The \passthrough{\lstinline!Lib/test/!} tree
includes broad tests for builtins, standard-library modules, interpreter
behavior and regression cases. Application repositories should test
boundaries where behavior is most likely to change: parsing, I/O,
persistence, native calls, concurrency and error handling.

\begin{lstlisting}[language=Python]
import unittest


def normalize_username(value: str) -> str:
    if not isinstance(value, str):
        raise TypeError("expected str")
    value = value.strip()
    if not 1 <= len(value) <= 40:
        raise ValueError("length out of range")
    return value.casefold()


class TestNormalizeUsername(unittest.TestCase):
    def test_success(self):
        self.assertEqual(normalize_username("  Alice  "), "alice")

    def test_missing(self):
        with self.assertRaises(ValueError):
            normalize_username("  ")

    def test_type(self):
        with self.assertRaises(TypeError):
            normalize_username(None)


if __name__ == "__main__":
    unittest.main()
\end{lstlisting}

An input-checking function demonstrates how to make a contract testable;
the specific username policy is \textbf{ENGINEERING}, not sourced from
CPython docs.

\section{Change impact and regression
testing}\label{change-impact-and-regression-testing}

For runtime changes, begin with the closest component test and expand
outward. A parser update requires syntax acceptance/rejection tests;
bytecode changes need semantic and exception tests; C API changes need
debug and multiple-build coverage; free-threading changes need
concurrency stress and lifetime verification. Investigate flaky tests by
controlling random seeds, resource constraints, process isolation and
timing rather than merely retrying until green.

\chapter{Security threat model for Python
services}\label{security-threat-model-for-python-services}

\section{Assets and entry points}\label{assets-and-entry-points}

Identify credentials, application data, logs, filesystem authority,
native libraries, subprocesses and network endpoints. Identify where
external data enters: HTTP requests, uploaded files, configuration, CLI
arguments, environment, deserialization streams, dependencies and
embedded interpreters. Model attacker control of types, sizes, encodings
and timing.

\textbf{Source emphasis:}
\path{Doc/library/security_warnings.rst} highlights
insecure uses of \passthrough{\lstinline!pickle!}, noncryptographic
\passthrough{\lstinline!random!}, limited
\passthrough{\lstinline!http.server!} security, logging configuration
evaluation, and security considerations for network and system modules.
These are specific warnings, not an exhaustive application threat model.

\section{High-impact unsafe patterns}\label{high-impact-unsafe-patterns}

\begin{longtable}[]{@{}
  >{\raggedright\arraybackslash}p{(\columnwidth - 4\tabcolsep) * \real{0.3333}}
  >{\raggedright\arraybackslash}p{(\columnwidth - 4\tabcolsep) * \real{0.3333}}
  >{\raggedright\arraybackslash}p{(\columnwidth - 4\tabcolsep) * \real{0.3333}}@{}}
\toprule\noalign{}
\begin{minipage}[b]{\linewidth}\raggedright
Trust boundary
\end{minipage} & \begin{minipage}[b]{\linewidth}\raggedright
Unsafe pattern
\end{minipage} & \begin{minipage}[b]{\linewidth}\raggedright
Safer direction
\end{minipage} \\
\midrule\noalign{}
\endhead
\bottomrule\noalign{}
\endlastfoot
Untrusted serialization &
\path{pickle.loads(attacker_data)} & Use JSON or
another non-executable format plus validation \\
Shell invocation & Interpolated command with
\passthrough{\lstinline!shell=True!} & Pass an argument vector; avoid
shell when possible \\
Credentials & \passthrough{\lstinline!random!} tokens &
\passthrough{\lstinline!secrets!} and suitable lifecycle policies \\
Code execution & \path{eval(user_input)} & Dedicated
restricted parser or data format \\
File paths & Joining untrusted path to privileged root unchecked &
Explicit authorization and race-aware file handling \\
Temporary files & Predictable filename and permissive mode &
\passthrough{\lstinline!tempfile!} secure creation patterns \\
TLS & Disabling certificate checking to fix errors & Fix
trust/configuration and fail closed \\
Logging & Dumping bearer tokens or full request bodies & Redact and
minimize \\
\end{longtable}

\begin{lstlisting}[language=Python]
import subprocess

def safe_probe(host: str) -> subprocess.CompletedProcess[str]:
    # Example only: validate the hostname and constrain execution further.
    if not host or host.startswith("-"):
        raise ValueError("invalid host")
    return subprocess.run(
        ["ping", "-c", "1", "--", host],
        text=True,
        capture_output=True,
        timeout=5,
        check=False,
    )
\end{lstlisting}

\textbf{Portability and safety note:} The \passthrough{\lstinline!ping!}
example uses a POSIX-style command and is not guaranteed to work on
every platform, and argument separation alone does not establish input
authorization or resource isolation. Define allowed host syntax, prevent
option injection and allowlist executables before using any command in a
privileged context.

\section{Deserialization safety}\label{deserialization-safety}

\path{Doc/library/pickle.rst} explicitly warns that
crafted pickle inputs may execute arbitrary code. HMAC authentication of
a pickle blob can establish tamper detection only when signing keys,
algorithms and replay policy are secure; it does not make
attacker-controlled creation of signed pickles safe. Prefer data formats
and schema validators that cannot instantiate arbitrary Python objects
during parsing.

\chapter{Static analysis and dependency
assurance}\label{static-analysis-and-dependency-assurance}

\textbf{ENGINEERING CI design:} run a pinned Python interpreter matrix,
syntax checks, linting, static type analysis appropriate to the project,
secret scanning, dependency vulnerability checks, unit tests, package
building, native extension verification, and deployment smoke tests.
Python annotations are not a runtime type enforcement mechanism; a type
checker supplements, rather than replaces, executable tests.

\textbf{Example GitHub Actions workflow (illustrative, not copied from
CPython):}

\begin{lstlisting}
name: python-assurance
on:
  push:
  pull_request:
jobs:
  verify:
    runs-on: ubuntu-latest
    strategy:
      matrix:
        python-version: ['3.13', '3.14']
    steps:
      - uses: actions/checkout@v4
      - uses: actions/setup-python@v5
        with:
          python-version: ${{ matrix.python-version }}
      - run: python -m compileall -q src tests
      - run: python -m unittest discover -s tests -v
\end{lstlisting}

Pin third-party actions to reviewed commit SHAs and package dependencies
to reproducible constraints in a production workflow. Version choices
are illustrations; test the versions actually supported by the
application. A development 3.16 alpha target should be an explicitly
opted-in compatibility job, not silently assumed stable.

\chapter{Performance engineering and
reproducibility}\label{performance-engineering-and-reproducibility}

\section{Define a meaningful
benchmark}\label{define-a-meaningful-benchmark}

A performance comparison needs a pinned workload, realistic input scale,
warm-up policy, repeat count, statistical summary, process and hardware
configuration, and a stated bottleneck hypothesis. Avoid benchmarking an
implementation micro-detail as if it were an API requirement.
\passthrough{\lstinline!timeit!} measures timed snippets;
\passthrough{\lstinline!cProfile!} helps attribute call costs;
\passthrough{\lstinline!tracemalloc!} reports Python-tracked
allocations. Native memory and system I/O may require extra tooling.

\begin{lstlisting}[language=Python]
from timeit import repeat

setup = "items = list(range(10000))"
times = repeat("sum(items)", setup=setup, repeat=7, number=500)
print("min seconds per 500 iterations:", min(times))
print("max seconds per 500 iterations:", max(times))
\end{lstlisting}

\section{Profiling decision process}\label{profiling-decision-process}

\begin{enumerate}
\def\labelenumi{\arabic{enumi}.}
\tightlist
\item
  Define the latency, throughput or memory SLO and percentile that
  matters.
\item
  Reproduce a production-like profile with representative input
  distributions.
\item
  Separate time in Python code, native extensions, I/O wait and lock
  contention.
\item
  Form one bottleneck hypothesis; optimize one cause at a time.
\item
  Re-run correctness and abuse tests; an optimization may alter cleanup
  or timing.
\item
  Report distributions, regressions, environmental configuration and
  measurement limits.
\end{enumerate}

\path{InternalDocs/jit.md} and
\path{InternalDocs/interpreter.md} describe
implementation strategies that can affect warm-up and opcode execution.
Interpretation of benchmark outcomes must therefore account for
build-time JIT/free-threading options and optimization state.

\chapter{Diagnostic workflow for production
failures}\label{diagnostic-workflow-for-production-failures}

\section{Symptoms to evidence}\label{symptoms-to-evidence}

\begin{longtable}[]{@{}
  >{\raggedright\arraybackslash}p{(\columnwidth - 4\tabcolsep) * \real{0.3333}}
  >{\raggedright\arraybackslash}p{(\columnwidth - 4\tabcolsep) * \real{0.3333}}
  >{\raggedright\arraybackslash}p{(\columnwidth - 4\tabcolsep) * \real{0.3333}}@{}}
\toprule\noalign{}
\begin{minipage}[b]{\linewidth}\raggedright
Symptom
\end{minipage} & \begin{minipage}[b]{\linewidth}\raggedright
First evidence
\end{minipage} & \begin{minipage}[b]{\linewidth}\raggedright
Likely follow-on investigation
\end{minipage} \\
\midrule\noalign{}
\endhead
\bottomrule\noalign{}
\endlastfoot
Exceptions & Traceback, error rate, correlation ID & Inputs, boundary
validation, regression \\
Long-tail latency & Metrics, timeout frequency & I/O timing, lock
contention, saturation \\
Memory growth & \passthrough{\lstinline!tracemalloc!}, RSS, object
counts & Retained objects, caches, native allocations \\
Concurrency hang & Thread/task dump, lock graph & Deadlock ordering,
missed cancellation \\
Native crash & Core dump, debug symbols, extension inventory & Reference
ownership, buffer safety, ABI \\
Different results by release & Version/platform/config metadata &
Language or implementation change \\
\end{longtable}

Do not log customer secrets or enable invasive tracing without an
access/retention plan. Diagnose fault containment before attempting an
in-place optimization.

\chapter{Release and deployment
checklist}\label{release-and-deployment-checklist}

\begin{enumerate}
\def\labelenumi{\arabic{enumi}.}
\tightlist
\item
  Capture interpreter version, supported platform matrix and dependency
  constraints.
\item
  Run all correctness tests and security abuse cases under at least the
  oldest and newest supported interpreter versions.
\item
  Verify secrets are supplied through approved runtime channels and not
  embedded in source or logs.
\item
  Inspect dependency licenses, provenance and vulnerabilities; decide on
  an explicit acceptance/deviation process.
\item
  Confirm packaging includes only intended files, signed artifacts where
  applicable and consistent build metadata.
\item
  Define runtime timeouts, memory limits, logging sensitivity and
  restoration procedure.
\item
  Exercise rollback, configuration failure, expired TLS certificates and
  service restart paths.
\item
  Document any implementation-specific assumptions (CPython internals,
  free-threading, GIL, C API).
\end{enumerate}

\chapter{Laboratories and assessment
cases}\label{laboratories-and-assessment-cases}

\textbf{Lab A --- deserialization risk:} write unit tests proving the
decoder rejects malformed or wrong-type JSON. Document why replacing it
with untrusted \passthrough{\lstinline!pickle.loads!} would create
code-execution risk. Do not construct exploit payloads to pass this lab.

\textbf{Lab B --- process invocation:} replace a string-building shell
invocation with an argument-vector design. Test rejected options,
timeouts, and error return codes.

\textbf{Lab C --- differential verification:} run a pure-Python API on
two pinned interpreter versions. Report behavior differences separately
from performance differences and cite the relevant documentation.

\textbf{Lab D --- memory investigation:} record
\passthrough{\lstinline!tracemalloc!} before/after snapshots for a
bounded workload and explain why Python-tracked allocation size can
differ from RSS.

\textbf{Lab E --- pipeline gate:} create a local CI-equivalent script
that runs syntax checks, \passthrough{\lstinline!unittest!} and
packaging validation. Intentionally introduce a failing test to show
that publishing does not proceed.

\textbf{Acceptance rubric:} 35\% correctness/test quality, 30\%
abuse-case completeness, 20\% observability/reproducibility, 15\%
remediation prioritization.

\chapter{Source concordance}\label{source-concordance}

\begin{itemize}
\tightlist
\item
  \textbf{SOURCE:}
  \path{Doc/library/security_warnings.rst},
  \passthrough{\lstinline!pickle.rst!},
  \passthrough{\lstinline!secrets.rst!},
  \path{subprocess.rst},
  \passthrough{\lstinline!tempfile.rst!},
  \passthrough{\lstinline!logging.rst!},
  \passthrough{\lstinline!timeit.rst!},
  \path{tracemalloc.rst},
  \path{unittest.mock-examples.rst}.
\item
  \textbf{SOURCE:}
  \path{InternalDocs/interpreter.md},
  \passthrough{\lstinline!jit.md!},
  \path{garbage_collector.md},
  \path{InternalDocs/README.md};
  \passthrough{\lstinline!Lib/test/!} repository structure.
\item
  \textbf{SOURCE:}
  \path{Doc/howto/free-threading-python.rst};
  \path{.github/workflows/} for an example of project
  CI structure, not specific third-party gate configurations.
\item
  \textbf{ENGINEERING:} threat model, CI workflow, test matrices,
  snippets and release checklist. Specific external security tools are
  illustrative extensions and are not claimed to be first-party CPython
  requirements.
\end{itemize}

\end{document}
